%! TEX root = main.tex

\subsection{zk-SNARK}
zk-SNARK(Succinct Non-interactive ARgument of Knowledge)는 범용 영지식 증명 기술 중 가장 대중적으로 널리 쓰이는 기술이다. 영지식 증명이란 어떠한 사실의 구체적인 내용은 밝히지 않으면서, 그것이 특정한 조건을 만족함을 수학적으로 증명하는 기술로, 데이터의 기밀성을 유지하면서도 무결성을 검증할 수 있게 한다. ‘범용(General-purpose)’ 영지식 증명이라 함은 특정 문제에 국한되지 않고, NP 복잡도 클래스에 속하는 모든 계산 문제를 증명할 수 있다는 의미이다.

\bigskip

zk-SNARK는 비 대화형 프로토콜(Non-interactive Protocol)의 일종이다. 이는 증명자와 검증자가 여러 번의 통신을 주고받아야 하는 대화형 프로토콜과 달리, 단 한 번의 증명 데이터 전송만으로 검증이 완료됨을 의미한다. 여기서 'Succinct(간결함)'는 증명 데이터의 크기가 매우 작고(수백 바이트 수준), 검증에 소요되는 시간이 증명하고자 하는 연산의 복잡도와 무관하게 거의 상수 시간(Constant time) 내에 이루어짐을 뜻한다.

\bigskip

zk-SNARK를 적용하기 위해서는 증명하고자 하는 연산(Computation)을 산술 회로(Arithmetic Circuit)의 형태로 표현해야 한다. 이 회로는 덧셈과 곱셈 게이트로 구성되며, 이는 다시 R1CS(Rank-1 Constraint System)라는 형태의 수학적 제약 조건 시스템으로 변환되어 처리된다. R1CS는 벡터의 내적 연산을 이용해 회로의 각 게이트가 올바르게 연산되었는지를 검증하는 방정식의 집합이다. 본 연구에서 활용하는 gnark 와 같은 프레임워크는 복잡한 R1CS 변환 과정을 추상화하여 제공함으로써, 개발자가 고수준의 코드로 회로를 정의할 수 있게 돕는다. 이러한 특성 덕분에 Zcash와 같은 블록체인에서의 트랜잭션 은닉 등 다양한 프라이버시 보호 시스템에서 핵심 기술로 활용되고 있다.

\subsection{Groth16 protocol}
Groth16 프로토콜은 현재 뛰어난 성능을 바탕으로 가장 널리 쓰이고 있는 zk-SNARK 프로토콜이다. 2016년 Jens Groth\cite{groth2016}가 처음 소개한 이 프로토콜은 증명 데이터의 크기와 검증 시간에 대해 상수 스케일의 복잡도를 달성하여 오늘날 zk-SNARK의 표준처럼 사용되고 있다. 본 연구 역시 해당 프로토콜을 기반으로 진행되었다.

\bigskip

증명자는 증명하고자 하는 선언(statement)을 다항식 세트 (Quadratic Arithmetic Program, QAP)로 표현한다. 예를 들어 증명자가 ‘$x^3+2x+7=19$를 만족하는 x값을 알고 있다’ 라는 사실(statement)를 x=2(witness)임을 밝히지 않으면서 증명을 한다면, 우선 x3+2x+7=19를 계산하는 과정을QAP라는 다항식 묶음의 형태로 치환해야 한다. 이 다항식들이 비밀 값(witness)을 내포하고 있는 다른 다항식으로 모두 나누어 떨어진다면, 검증자는 증명자가 옳은 비밀 값으로 그 QAP를 구성했다는 확신을 얻을 수 있다. 이 과정은 두 다항식이 하나의 임의로 선택된 점에서 교차한다면 충분히 높은 확률로 그 두 다항식이 동일하다는 Schwartz—Zippel Lemma에 의해 빠르고 효율적인 검증임이 증명된다. 그리고 암호화된 상태에서의 검증을 위해 이산로그 기반의 타원곡선 암호화 알고리즘, bilinear pairing등 여러 암호학적 장치들을 통해 이 검증 과정에서 증명 사실 외 다른 정보를 알 수 없도록 숨긴다.

\bigskip

Groth16은 비 대화형 프로토콜을 구현하기 위해 신뢰 셋업(trusted setup)이라는 공유 파라미터들을 증명 이전에 생성한다. 신뢰 셋업이란 검증자가 비밀의 내용을 모르고도 선언의 내용을 검증할 수 있게 하는 암호학적 파라미터들을 말한다. 대화형 프로토콜에선 증명자와 검증자가 서로 임의의 값들을 주고 받으며 처리해야 할 계산들을 이렇게 미리 합의된 값을 이용해 한번에 처리하는 방식이다. 검증자는 증명자가 제시한 다항식 위에서 임의의 한 점을 선택해 그에 대한 계산 결과를 검증해야 한다. 그런데 이 점이 증명자에게 노출되면 증명자는 해당 점을 포함하도록 다항식을 조작할 수 있다. 따라서 신뢰 셋업은 공개적인 과정을 통해 파라미터를 생성하여 그 임의의 점에 대해 서로 합의하는 과정이라고 볼 수 있다. 이 과정에서 Powers of Tau라 불리는 MPC(Multi-Party Computation) 프로토콜이 사용되며, 이는 참여자가 각자가 선택한 임의의 비밀 값 ��(tau)를 엔트로피로써 파라미터에 추가하는 과정으로 요약할 수 있다. 이 참여자 중 단 한명이라도 ��를 제대로 폐기했다면 이 신뢰 셋업의 기밀성은 수학적으로 보장되는 방식이다.

\subsection{Merkle Proof}
머클 트리(Merkle Tree)는 1979년 Ralph C. Merkle이 제안한 인증된 데이터 구조(Authenticated Data Structure, ADS)의 일종이다. 이는 대량의 데이터셋을 해시 함수를 이용해 이진 트리 형태로 요약한 구조로, 전체 데이터를 노출하지 않고도 특정 데이터가 해당 셋에 포함되어 있음을 효율적으로 증명(Membership Proof)하는 데 사용된다.

\bigskip

머클 트리의 핵심은 루트 노드(Merkle Root)가 트리 내의 모든 데이터에 대한 암호학적 압축값(Fingerprint) 역할을 한다는 점이다. 전체 데이터의 개수가 n일 때, 특정 리프 노드가 트리에 포함되어 있음을 증명하기 위해 전체 데이터를 검증할 필요 없이, 루트까지의 경로에 있는 해시값들(Sibling hashes)만 있으면 되므로 검증 복잡도는 O(log n)으로 매우 효율적이다.

\bigskip

본 연구는 이 머클 증명을 데이터 발급자(Data Issuer)의 유효성 검증에 활용한다. 시스템에 등록된 수많은 발급자의 공개키 중, 현재 데이터를 서명한 키가 유효한 목록에 포함되어 있다는 사실을 증명해야 한다. 이때 머클 증명을 zk-SNARK 회로 내부에서 수행함으로써, 검증자에게는 '유효한 발급자 중 한 명'이라는 사실만 알리고, 구체적으로 '누구'인지는 밝히지 않는 익명성을 보장한다. 이처럼 머클 증명은 블록체인의 트랜잭션 검증뿐만 아니라, 프라이버시를 보존하는 멤버십 증명 도구로서도 중요한 역할을 수행한다.

\subsection{zk-SNARK 친화적 암호화 원형 (zk-SNARK Friendly Cryptographic Primitives)}
zk-SNARK 시스템의 성능은 증명하고자 하는 연산을 산술 회로(Arithmetic Circuit)로 변환했을 때 생성되는 제약 조건(Constraints)의 수에 크게 의존한다. 일반적인 암호화 알고리즘(예: SHA-256, RSA)은 비트 연산이 많아 산술 회로로 구현 시 제약 조건의 수가 폭발적으로 증가하여 증명 생성 시간을 지연시킨다. 따라서 본 연구에서는 회로 내부에서의 연산 비용을 최소화하기 위해 대수적 구조가 간단한 zk-SNARK 친화적(zk-friendly) 암호화 원형들을 채택하여 사용한다.

\subsubsection{EdDSA (Edwards-curve Digital Signature Algorithm)} 
본 시스템은 데이터 발급자(DI)의 신원을 회로 내부에서 검증하기 위해 EdDSA 서명 알고리즘을 사용한다. EdDSA는 타원곡선 암호(ECC)의 일종으로, 특히 Twisted Edwards Curve 위에서 정의된다. 본 연구에서는 zk-SNARK 회로(주로 BN254 또는 BLS12-381 곡선 위에서 작동) 내부에 임베딩하기 적합한 Baby Jubjub 타원곡선을 사용한다. 이 곡선은 회로의 기본 필드(Base Field)와 호환되도록 설계되어 있어, 서명 검증 시 발생하는 타원곡선 덧셈 및 스칼라 곱셈 연산을 매우 적은 수의 제약 조건만으로 효율적으로 수행할 수 있다. 이는 수천 번의 곱셈이 필요한 RSA나 ECDSA에 비해 압도적인 성능 우위를 제공한다.

\subsubsection{MiMC (Minimal Multiplicative Complexity) Hash} 
데이터의 무결성을 검증하고 머클 트리(Merkle Tree)를 구성하기 위해 본 연구는 MiMC 해시 함수를 사용한다. 일반적으로 사용되는 SHA-256 해시 함수는 비트 단위의 논리 연산(XOR, Shift)을 다량으로 사용하므로 산술 회로 구현 시 수만 개의 제약 조건을 유발한다. 반면, MiMC는 유한체(Finite Field) 상에서의 세제곱 연산(x3)이나 역수 연산(x−1)과 같은 간단한 대수적 변환을 반복하는 구조로 설계되어 있다. 이러한 구조는 zk-SNARK의 산술 회로(R1CS) 표현에 최적화되어 있어, SHA-256 대비 수십 배 적은 제약 조건으로 동일한 보안 강도의 해시 연산을 수행할 수 있게 한다. 본 시스템의 머클 경로 검증(Merkle Path Verification) 로직은 MiMC를 사용하여 회로 깊이를 획기적으로 줄였다.

\subsubsection{암호학적 커밋 (Cryptographic Commitment)} 
데이터의 기밀성을 보장하면서도 특정 데이터가 변경되지 않았음을 증명하기 위해 암호학적 커밋 기법이 활용된다. 본 시스템에서는 다수의 검색된 증명 데이터들이 동일한 소유자(DP)의 것임을 영지식으로 증명하는 과정에서 정보를 노출하지 않기 위해 해시 함수 기반의 커밋을 사용한다. 특히 zk-SNARK 회로 내부에서의 연산 효율성을 극대화하기 위해 앞서 언급한 MiMC 해시 함수를 이용하여 소유자의 식별자(DP\_ID)와 증명 식별자들을 결합한 커밋을 생성하며, 이를 통해 완전한 은닉성(Hiding)과 위조 불가능성(Binding)을 달성한다.

\bigskip

이와 같이 본 연구는 회로 연산 비용에 최적화된 암호화 원형들을 전략적으로 선택 및 조합하여, 높은 보안성을 유지하면서도 실용적인 증명 생성 및 검증 속도를 달성한다.
